\title{DeepSeekMath: Pushing the Limits of Mathematical Reasoning in Open Language Models}

\begin{abstract}
Language models frequently encounter substantial difficulties in mathematical reasoning because of its inherent complexity and structured requirements. In this study, we present DeepSeekMath~7B, developed by further pre-training the DeepSeek-Coder-Base-v1.5 7B on a corpus containing 120B math-centric tokens, drawn from Common Crawl and supplemented with natural language and programming code data. DeepSeekMath~7B attains an outstanding result of 51.7\% on the competition-level MATH benchmark, all without the aid of external toolkits or ensemble-based voting approaches, and demonstrates capabilities on par with Gemini-Ultra and GPT-4. When evaluated for self-consistency across 64 generated samples, DeepSeekMath~7B reaches 60.9\% accuracy on the MATH benchmark. The enhanced mathematical reasoning found in DeepSeekMath~7B is primarily enabled by two aspects: first, leveraging a carefully crafted data selection pipeline to exploit the breadth of accessible web data; and second, the introduction of Group Relative Policy Optimization (GRPO), an adaptation of Proximal Policy Optimization (PPO), which simultaneously augments mathematical reasoning skills and improves memory efficiency within PPO.
\end{abstract}

\section{Introduction}

Large language models (LLM) have transformed the field of mathematical reasoning within artificial intelligence, driving notable progress in benchmarks like quantitative reasoning \citep{MATH} and geometry reasoning \citep{trinh2024solving}. In addition, these models serve as valuable tools for supporting human efforts in addressing sophisticated mathematical challenges \citep{tao}. Despite these achievements, leading models such as GPT-4 \citep{gpt4} and Gemini-Ultra \citep{gemini} remain proprietary, while open-source alternatives that are available demonstrate a marked gap in performance by comparison.

This paper presents DeepSeekMath, a specialized language model tailored for mathematical tasks, which notably surpasses the mathematical reasoning abilities of other open-source systems and performs comparably to GPT-4 on a range of academic benchmarks. Central to this achievement is the construction of the DeepSeekMath Corpus—a comprehensive, high-caliber pre-training dataset containing 120B math tokens. Data for this corpus are sourced from Common Crawl (CC) through the utilization of a fastText-based classifier \citep{joulin2016fasttext}. The classifier, in its initial phase, is trained with positive samples taken from OpenWebMath \citep{openwebmath} and a heterogeneous mix of general web content as negative samples. Next, the classifier is used to extract further positive examples from the CC dataset; these instances undergo additional refinement through human annotation. The classifier is subsequently retrained with these enhanced labels, leading to improved extraction performance. Empirical results reveal that the resultant large-scale dataset exhibits exceptional quality: the DeepSeekMath-Base 7B model records 64.2\% accuracy on GSM8K \citep{gsm8k} and 36.2\% on the high-difficulty MATH benchmark \citep{MATH}, achieving better results than Minerva 540B \citep{minerva}. Furthermore, thanks to the multilingual design of the DeepSeekMath Corpus, we observe notable progress on Chinese-language mathematical benchmarks \citep{wei2023cmath,agieval}. We contend that our systematic approach to curating mathematical data can serve as a valuable reference point for the field, with considerable opportunities for future refinement and enhancement.

DeepSeekMath-Base is built upon DeepSeek-Coder-Base-v1.5 7B \citep{deepseek-coder}, as initiating the process with a code-specialized pre-trained model proves more effective than beginning with a standard LLM. Additionally, our results show that training with mathematical data not only boosts mathematical performance but also leads to improvements on the MMLU \citep{mmlu} and BBH \citep{bbh} benchmarks. This suggests that such training enhances the model's overall reasoning skills, extending beyond just mathematics.

Following pre-training, we conduct mathematical instruction tuning on DeepSeekMath-Base utilizing data that includes chain-of-thought \citep{cot}, program-of-thought \citep{pot,pal}, and tool-integrated reasoning \citep{tora}. This yields DeepSeekMath-Instruct 7B, which surpasses all other 7B models in its class and demonstrates performance on par with leading 70B open-source instruction-tuned models.

In addition, we present Group Relative Policy Optimization (GRPO), an alternative reinforcement learning (RL) approach derived from Proximal Policy Optimization (PPO) \citep{schulman2017proximal}. Unlike PPO, GRPO dispenses with the critic model, instead leveraging group scores as baseline estimates, which greatly lowers the computational cost of training. Leveraging only a portion of English instruction tuning data, GRPO achieves notable gains over the already strong DeepSeekMath-Instruct, showing marked improvements on in-domain benchmarks (GSM8K: 82.9\% $\rightarrow$ 88.2\%, MATH: 46.8\% $\rightarrow$ 51.7\%) as well as enhanced performance on out-of-domain math tasks (e.g., CMATH: 84.6\% $\rightarrow$ 88.8\%) during the reinforcement learning phase. We further offer a unified framework to interpret various methodologies, including Rejection Sampling Fine-Tuning (RFT) \citep{yuan2023scaling}, Direct Preference Optimization (DPO) \citep{dpo}, PPO, and GRPO. This unified perspective reveals that these strategies can all be understood as forms of direct or streamlined RL algorithms. We carry out thorough experimentation—such as comparing online and offline training, outcome-based versus process-based supervision, and single-step versus iterative RL—to systematically analyze the fundamental aspects of this paradigm. Finally, we elucidate how our RL approach enhances the capabilities of instruction-tuned models, and outline promising avenues for future improvements to RL methods within this framework.

\subsection{Contributions}
We present advancements in scalable math pre-training, as well as investigate and assess the use of reinforcement learning.

\textbf{Math Pre-Training at Scale}

\begin{itemize}[topsep=0pt]
    \item Through our investigation, we demonstrate that the openly available Common Crawl resource harbors substantial mathematical content of high utility. Utilizing a carefully constructed data curation pipeline, we developed the DeepSeekMath Corpus—a large-scale, high-quality dataset comprising 120B tokens extracted from web pages with a focus on mathematical relevance. This corpus is nearly seven times larger than the math-specific subset employed by Minerva \citep{minerva}, and approximately nine times greater in size compared to the recently published OpenWebMath \citep{openwebmath}.

\item DeepSeekMath-Base 7B, our base model, attains results on par with Minerva 540B \citep{minerva}. This suggests that sheer parameter count is not solely responsible for proficiency in mathematical reasoning. Smaller models, when trained on curated, high-quality datasets, can also deliver impressive outcomes.

\item We present the results from our math training experiments. Pre-training models on code before introducing math training enhances their performance on mathematical tasks, regardless of whether they use external tools. This contributes to the ongoing debate regarding the impact of code training on reasoning: \textit{does code training improve reasoning abilities?} Our results suggest that it does, particularly in the context of mathematical reasoning.

\item While it is standard practice to utilize arXiv papers during training, particularly within the mathematics domain, this approach does not yield any significant performance gains across the mathematical benchmarks used in this study.
\end{itemize}

\textbf{Exploration and Analysis of Reinforcement Learning}

\begin{itemize}[topsep=0pt]
    \item We propose Group Relative Policy Optimization (GRPO), a novel reinforcement learning approach characterized by its efficiency and effectiveness. Unlike traditional methods that rely on a critic network, GRPO leverages group scores to estimate baselines, thus markedly decreasing computational demands compared to Proximal Policy Optimization (PPO).
    \item We provide empirical evidence that GRPO markedly improves the capabilities of our instruction-tuned model, DeepSeekMath-Instruct, utilizing only instruction-tuning datasets. In addition, the reinforcement learning phase leads to noticeable gains in generalization performance for out-of-domain tasks.
    \item We present a comprehensive framework that encompasses various methodologies, including RFT, DPO, PPO, and GRPO. Furthermore, we perform thorough experimental evaluations—such as comparing online and offline learning, outcome-based versus process-based supervision, and single-step versus iterative reinforcement learning—to systematically explore key components of this framework.
    \item Drawing on this unified framework, we examine the underlying factors contributing to reinforcement learning's efficacy and outline several promising research directions for further enhancing the reinforcement learning capabilities of large language models (LLMs).
\end{itemize}

\subsection{Summary of Evaluations and Metrics}

\begin{itemize}[topsep=0pt]
    \item \textbf{English and Chinese Mathematical Reasoning}: Our evaluation extensively analyzes model performance across both English and Chinese mathematical datasets, spanning questions from elementary through collegiate levels. For English, we utilize benchmarks such as GSM8K \citep{gsm8k}, MATH \citep{MATH}, SAT \citep{llemma}, OCW Courses \citep{minerva}, and MMLU-STEM \citep{mmlu}. Chinese assessment includes MGSM-zh \citep{mgsm}, CMATH \citep{wei2023cmath}, Gaokao-MathCloze \citep{agieval}, and Gaokao-MathQA \citep{agieval}. Our analysis examines model proficiency in composing independent written solutions and in addressing problems with Python-based tool support.

On English benchmark tasks, DeepSeekMath-Base matches the performance of the proprietary Minerva 540B model \citep{minerva} and outperforms all open-source foundational models, such as Mistral 7B \citep{mistral} and Llemma-34B \citep{llemma}, irrespective of whether these models have been pre-trained on mathematical data, often by a noteworthy margin. Importantly, DeepSeekMath-Base excels on Chinese benchmark evaluations, which may be attributed to our decision not to limit mathematical pre-training data to English as in prior studies \citep{minerva,llemma}, but to include carefully selected high-quality non-English resources as well. Through the application of mathematical instruction tuning and reinforcement learning, the derived models DeepSeekMath-Instruct and DeepSeekMath-RL achieve robust results, surpassing the 50\% accuracy threshold on the competition-level MATH dataset—marking a first among open-source efforts.

\item \textbf{Formal Mathematics}: DeepSeekMath-Base is assessed via the informal-to-formal theorem proving benchmark introduced in \citep{dsp_proof}, utilizing miniF2F \citep{minif2f} and employing Isabelle \citep{isabelle} as the designated proof assistant. The results indicate that DeepSeekMath-Base achieves robust few-shot autoformalization outcomes.
\item \textbf{Natural Language Understanding, Reasoning, and Code}: To thoroughly evaluate the general proficiency of models in understanding, reasoning, and programming, DeepSeekMath-Base is benchmarked on the Massive Multitask Language Understanding (MMLU) suite \citep{mmlu}, which spans 57 multiple-choice tasks from a broad spectrum of fields; BIG-Bench Hard (BBH) \citep{bbh}, containing 23 difficult tasks that predominantly entail complex multi-step reasoning; as well as HumanEval \citep{codex} and MBPP \citep{mbpp}, both standard benchmarks for coding capabilities in language models. Pre-training on mathematical data contributes positively to both language comprehension and reasoning performance.

\section{Math Pre-Training}

\subsection{Data Collection and Decontamination} 

This section details the methodology used to assemble the DeepSeekMath Corpus utilizing data from Common Crawl. Figure \ref{fig:data_collect} illustrates an iterative workflow employed to efficiently extract a substantial mathematical dataset from Common Crawl, initially leveraging a seed corpus—typically a concise yet high-quality math-centric dataset. Importantly, this strategy can be generalized and applied across various domains, including coding.

Initially, we select the OpenWebMath dataset \citep{openwebmath}, which comprises a diverse array of well-curated mathematical texts from the web, as our foundational seed dataset. Using this seed, we construct a fastText model \citep{joulin2016fasttext} for the retrieval of additional web pages with characteristics similar to those in OpenWebMath. For model training, we randomly draw 500,000 positive examples from our seed dataset and 500,000 negative examples from the Common Crawl web archive. The training utilizes an open-source implementation\footnote{\url{https://fasttext.cc}}, setting the vector dimensionality to 256, the learning rate at 0.1, a maximum n-gram length of three, a minimum word occurrence threshold of three, and three training epochs. To reduce redundancy in Common Crawl, we perform URL and near-duplicate filtering, yielding a processed collection of 40B HTML pages. We apply the trained fastText classifier to this deduplicated corpus, retrieving those web pages most likely to contain mathematical content. Next, we filter and prioritize these pages using their respective fastText scores, retaining only the highest ranked for further use. To determine the data subset suitable for pre-training, we evaluate the preserved corpus by conducting experiments on subsets of 40B, 80B, 120B, and 160B tokens. For the first iteration, we opt to preserve the top 40B tokens.

Following the initial round of data gathering, a substantial portion of mathematical web pages remain uncollected. This is primarily due to the fastText model being trained on a positive example set that does not exhibit adequate diversity. To enhance the fastText model's performance, we expand the seed corpus by identifying more sources of mathematical web content. Our strategy begins by partitioning the entire Common Crawl dataset into non-overlapping domains, where each domain represents all web pages sharing a common base URL. For each domain, we determine the fraction of pages gathered during the first iteration. Domains in which more than 10\% of pages are collected are designated as math-related (such as \url{mathoverflow.net}). 

We then manually review and annotate the URLs pointing to mathematical content within these math-related domains (for instance, \url{mathoverflow.net/questions}). Web pages linked to these annotated URLs that were not captured previously are incorporated into the seed corpus. Through this process, we augment our collection of positive examples, which facilitates the training of a more robust fastText model with improved recall for mathematical web content in subsequent iterations. By the end of four iterative cycles, we accumulate a dataset comprising 35.5M mathematical web pages, amounting to 120B tokens in total. During the fourth iteration, we observe that approximately 98\% of the pages were already obtained in the previous round, prompting us to terminate further data collection efforts.

To prevent benchmark contamination, we adopt the filtering procedure outlined in \cite{deepseek-coder}, removing web content that includes questions or answers from English mathematical benchmarks like GSM8K \citep{gsm8k} and MATH \citep{MATH}, as well as Chinese benchmarks such as CMATH \citep{wei2023cmath} and AGIEval \citep{agieval}. Specifically, any segment of text that contains a 10-gram sequence precisely matching a substring from any of the evaluation benchmarks is excluded from our mathematics training dataset. For benchmark texts consisting of fewer than 10 grams but no less than 3 grams, we apply strict exact matching to identify and eliminate potentially contaminated web pages.

\subsection{Validating the Quality of the DeepSeekMath~Corpus}

We run pre-training experiments to investigate how the DeepSeekMath~Corpus is compared with the recently released math-training corpora:

\begin{itemize}[topsep=0pt]
    \item \textbf{MathPile} \citep{mathpile}: This collection combines 8.9B tokens from various sources, including textbooks, Wikipedia, ProofWiki, CommonCrawl, StackExchange, and arXiv, with arXiv constituting more than 85\% of the total data;
    \item \textbf{OpenWebMath} \citep{openwebmath}: Extracted from CommonCrawl by selecting for mathematical material, this dataset comprises 13.6B tokens;
    \item \textbf{Proof-Pile-2} \citep{llemma}: This dataset merges OpenWebMath, AlgebraicStack (with 10.3B tokens focusing on mathematical programming), and 28.0B tokens from arXiv manuscripts. In our experiments with Proof-Pile-2, we employ the arXiv:Web:Code proportion of 2:4:1, following the procedure in \cite{llemma}.
\end{itemize}

\subsubsection{Training Setting}
\label{sec:quality-policy}
Mathematics-specific training is conducted on a foundational language model consisting of 1.3B parameters, built upon the same architecture as DeepSeek LLMs \citep{deepseek-llm}, referred to here as DeepSeek-LLM 1.3B. The model is trained independently on each mathematical dataset for a total of 150B tokens. All training procedures utilize the HAI-LLM framework \citep{haillm}, which is designed for efficiency and scalability. In alignment with the approaches used for DeepSeek LLMs, the AdamW optimizer \citep{adamW} is employed, configured with $\beta_1=0.9$, $\beta_2=0.95$, and $\mathrm{weight\_decay}=0.1$. The learning rate follows a multi-phase schedule: it ramps up to its maximum value after 2{,}000 warm-up steps, drops to 31.6\% of its peak at 80\% completion of training, and further tapers to 10.0\% at 90\% progress. The peak learning rate is set to 5.3e-4. Each training batch processes 4M tokens, maintaining a context window of 4K tokens.

\subsubsection{Evaluation Results}

\textbf{The DeepSeekMath~Corpus is of high quality, covers multilingual mathematical content, and is the largest in size.}

\begin{itemize}[topsep=0pt]
    \item \textbf{High-quality}: 
    We assess model performance on eight mathematical benchmarks using few-shot chain-of-thought prompting as described in \cite{cot}. Table \ref{tab:corpora_comparison} reveals a distinct advantage for models trained with the DeepSeekMath~Corpus, which consistently surpass their counterparts. Furthermore, Figure \ref{fig:corpora_comparisons} illustrates that the DeepSeekMath Corpus-trained model outperforms the Proof-Pile-2-trained counterpart at 50B tokens (corresponding to one complete pass over Proof-Pile-2). This suggests that the DeepSeekMath Corpus is of a generally superior quality.
    \item \textbf{Multilingual}: 
    The DeepSeekMath~Corpus is comprised of resources in various languages, with English and Chinese being the most prevalent. According to Table \ref{tab:corpora_comparison}, utilizing the DeepSeekMath~Corpus for training leads to improved mathematical reasoning in both English and Chinese. By contrast, the predominantly English-based existing corpora offer marginal gains and can even be detrimental to mathematical reasoning performance in Chinese.
    \item \textbf{Large-scale}:
    The DeepSeekMath~Corpus significantly exceeds previous mathematical corpora in size. As illustrated in Figure \ref{fig:corpora_comparisons}, DeepSeek-LLM 1.3B, when trained on this dataset, exhibits a more rapid and sustained enhancement in performance. In comparison, models trained on smaller, repeatedly cycled baseline corpora quickly experience stagnation in performance.
\end{itemize}

\subsection{Training and Evaluating DeepSeekMath-Base 7B}

This section presents DeepSeekMath-Base 7B, a foundational model that exhibits robust mathematical reasoning skills. The initial weights are based on DeepSeek-Coder-Base-v1.5 7B \citep{deepseek-coder}, after which the model is further pretrained using a dataset of 500B tokens. The composition of the training data is as follows: 56\% derives from the DeepSeekMath Corpus, 4\% from AlgebraicStack, 10\% from arXiv, 20\% consists of Github code, and the final 10\% comprises natural language content extracted from Common Crawl in both English and Chinese. The model is largely trained under the configuration detailed in Section \ref{sec:quality-policy}, but with a maximum learning rate of 4.2e-4 and a batch size of 10M tokens.

We carry out an in-depth evaluation of the mathematical proficiency exhibited by DeepSeekMath-Base 7B, with an emphasis on its competence in generating standalone mathematical solutions independently of auxiliary tools, employing tools for solving mathematical tasks, and performing formal proofs of theorems. In addition to its mathematical performance, we present a broader overview of the base model’s general capabilities, covering aspects such as natural language comprehension, logical reasoning, and programming abilities.

\paragraph{Mathematical Problem Solving with Step-by-Step Reasoning}

We assess the capabilities of DeepSeekMath-Base in addressing mathematical problems by utilizing few-shot chain-of-thought prompting \citep{cot}, across a set of eight benchmarks presented in both English and Chinese. These evaluation sets include tasks on quantitative reasoning—such as GSM8K \citep{gsm8k}, MATH \citep{MATH}, and CMATH \citep{wei2023cmath}—as well as multiple-choice assessments like MMLU-STEM \citep{mmlu} and Gaokao-MathQA \citep{agieval}. The benchmarks collectively span various mathematical domains, encompassing a spectrum of complexity from elementary through collegiate levels.

Table \ref{tab:base_math_cot} demonstrates that DeepSeekMath-Base 7B consistently achieves the highest scores across all eight evaluated benchmarks among open-source base models. This comparison includes prominent models such as the versatile Mistral 7B \citep{mistral} and the newly introduced Llemma 34B \citep{llemma}, which received additional mathematical training using the Proof-Pile-2 dataset \citep{llemma}. Particularly impressive is DeepSeekMath-Base's performance on the challenging MATH competition dataset, where it exceeds the accuracy of all open-source base models by a margin of over 10 percentage points. Furthermore, it outperforms the much larger Minerva 540B model \citep{minerva}—a proprietary base model based on PaLM \citep{palm} and subjected to extensive mathematical text training—even though Minerva contains 77 times more parameters.

\paragraph{Mathematical Problem Solving with Tool Use} We assess the capability of models in mathematical reasoning supported by tool use on the GSM8K and MATH benchmarks by employing few-shot program-of-thought prompting \citep{pot,pal}. The methodology involves instructing models to address each question by generating Python code that can access libraries like \textit{math} and \textit{sympy} to handle complex calculations. The computed output from the executed code serves as the final answer. As indicated in Table \ref{tab:base_math_pot_proof}, DeepSeekMath-Base 7B achieves superior results compared to the previously best-performing Llemma 34B.

\paragraph{Formal Mathematics}

The automation of formal proofs offers significant advantages in guaranteeing the correctness and robustness of mathematical demonstrations, while also improving productivity. This area has seen a surge of interest in recent years. In this work, we assess the DeepSeekMath-Base 7B model for the informal-to-formal proof generation task as described in \citep{dsp_proof}: generating a formal proof given an informal problem statement, its formal representation, and an associated informal proof. Our evaluation uses the miniF2F dataset \citep{minif2f}, a benchmark designed for formalizing mathematics at the Olympiad level, where we prompt the model with few-shot examples to produce formal proofs in Isabelle for each problem. Adopting the methodology of \cite{dsp_proof}, we have the models output proof sketches, which are then completed using Sledgehammer \citep{sledgehammer}, an automated theorem prover. Table \ref{tab:base_math_pot_proof} reports that DeepSeekMath-Base 7B achieves notable results in the task of automated proof formalization.

\paragraph{Natural Language Understanding, Reasoning, and Code}

We assess the capabilities of the model in natural language comprehension using MMLU \citep{mmlu}, in reasoning with BBH \citep{bbh}, and in programming tasks on HumanEval \citep{codex} and MBPP \citep{mbpp}. Table \ref{tab:understanding_reasoning_code} demonstrates that DeepSeekMath-Base 7B achieves marked improvements in both MMLU and BBH scores when compared with its predecessor, DeepSeek-Coder-Base-v1.5 \citep{deepseek-coder}, highlighting the beneficial effects of mathematical training on understanding and reasoning tasks. Furthermore, thanks to the integration of code tokens during continual learning, DeepSeekMath-Base 7B is able to retain the coding performance level of DeepSeek-Coder-Base-v1.5 across the evaluated code benchmarks. In summary, DeepSeekMath-Base 7B surpasses the general-purpose Mistral 7B model \citep{mistral} by a substantial margin on all three assessed reasoning and coding tasks.

\section{Supervised Fine-Tuning}

\subsection{SFT Data Curation}

We develop an instruction-tuning dataset for mathematics that encompasses both English and Chinese questions, spanning diverse mathematical domains and complexity grades. Each question is accompanied by a solution in one of several formats: chain-of-thought (CoT) \citep{cot}, program-of-thought (PoT) \citep{pot,pal}, or a format integrating external tools \citep{tora}. Altogether, the dataset comprises 776K training samples.

\begin{itemize}[topsep=0pt]
    \item \textbf{English mathematical datasets}: We provide tool-supported annotations for GSM8K and MATH problems, and utilize a selected portion of MathInstruct \citep{MathInstruct} as well as the Lila-OOD training set \citep{lila}, where solutions are presented using either CoT or PoT approaches. The assembled English dataset encompasses a wide array of mathematical disciplines, including but not limited to algebra, probability, number theory, calculus, and geometry.
    \item \textbf{Chinese mathematical datasets}: We assemble a collection of Chinese K-12 math problems across 76 distinct sub-topics, such as linear equations, with each problem annotated in both CoT and tool-based reasoning formats.
\end{itemize}

\subsection{Training and Evaluating DeepSeekMath-Instruct 7B}

In this section, we present DeepSeekMath-Instruct 7B, which is built upon DeepSeekMath-Base and refined through instruction tuning with mathematical tasks. Training data samples are assembled in random order and combined until the total reaches a 4K token context length. The model is trained over 500 steps, utilizing a batch size of 256 and maintaining a fixed learning rate of $5\times10^{-5}$.

We evaluate models' mathematical performance both without and with tool use, on 4 quantitative reasoning benchmarks in English and Chinese. We benchmark our model against the leading models of the time:

\begin{itemize}[topsep=0pt]
    \item \textbf{Closed-source models} comprise: (1) the GPT series, with GPT-4 \citep{gpt4} and GPT-4 Code Interpreter~\footnote{\url{https://openai.com/blog/chatgpt-plugins\#\#code-interpreter}} representing the highest-performing variants, (2) Gemini Ultra and Pro \citep{gemini}, (3) Inflection-2 \citep{inflection-2}, (4) Grok-1~\footnote{\url{https://x.ai/model-card}}, in addition to a selection of more recently introduced models from Chinese enterprises, including (5) Baichuan-3~\footnote{\url{https://www.baichuan-ai.com}} and (6) the most recent GLM-4~\footnote{\url{https://open.bigmodel.cn/dev/api\#glm-4}}

    from the GLM family \citep{glm}.

These models are designed for broad applicability, and the majority have been subjected to various alignment strategies. 
    \item \textbf{Open-source models} in this context encompass: general-purpose systems such as (1) DeepSeek-LLM-Chat 67B \citep{deepseek-llm}, (2) Qwen 72B \citep{qwen}, (3) SeaLLM-v2 7B \citep{seallm}, and (4) ChatGLM3 6B \citep{chatglm3}. Also included are models oriented towards mathematical capability, specifically: (5) InternLM2-Math 20B~\footnote{\url{https://github.com/InternLM/InternLM-Math}}, which extends InternLM2 and undergoes both targeted mathematical pre-training and subsequent instruction-based fine-tuning; (6) Math-Shepherd-Mistral 7B, which leverages PPO-based optimization \citep{schulman2017proximal} of Mistral 7B \citep{mistral} utilizing a reward model supervised on problem-solving processes; (7) WizardMath models \citep{wizardmath}, which augment the mathematical reasoning ability of Mistral 7B and Llama-2 70B \citep{llama2} using an evolve-instruct pipeline—an adaptation of instruction tuning wherein instructions are autonomously generated by AI—and reinforcement learning with PPO on tasks drawn largely from the GSM8K and MATH datasets; (8) MetaMath 70B \citep{metamath}, a version of Llama-2 70B that has been further adapted using extended sets of GSM8K and MATH; (9) ToRA 34B \cite{tora}, in which CodeLlama 34B is specialized through fine-tuning for tool-augmented mathematical reasoning; and (10) MAmmoTH 70B \citep{MathInstruct}, which relies on instruction-tuned Llama-2 70B using the MathInstruct dataset.

As indicated in Table \ref{tab:sft_rl_math}, when tool usage is prohibited during evaluation, DeepSeekMath-Instruct 7B excels in executing step-by-step mathematical reasoning. Remarkably, on the challenging MATH competition dataset, this model outperforms all open-source alternatives and the vast majority of commercial systems (including Inflection-2 and Gemini Pro) by a margin of at least 9\% in absolute terms. This performance advantage holds even over significantly larger models (such as Qwen 72B) and those that have undergone dedicated reinforcement learning aimed at mathematical capabilities (for instance, WizardMath-v1.1 7B). Although DeepSeekMath-Instruct is competitive with leading Chinese commercial models like GLM-4 and Baichuan-3 on the MATH benchmark, its results still fall short compared to GPT-4 and Gemini Ultra.

In the evaluation scenario permitting models to employ both natural language reasoning and programmatic tool usage for solving problems, DeepSeekMath-Instruct 7B achieves nearly 60\% accuracy on the MATH benchmark, outperforming all previously available open-source models. Furthermore, across additional benchmarks, our model demonstrates performance on par with DeepSeek-LLM-Chat 67B—the former state-of-the-art model that is ten times larger in scale.
\section{Reinforcement Learning}

\subsection{Group Relative Policy Optimization} Prior research has demonstrated that reinforcement learning (RL) significantly enhances the mathematical reasoning skills of LLMs beyond what is achieved through Supervised Fine-Tuning (SFT) alone \citep{wang2023math,wizardmath}. Here, we present our RL approach, termed Group Relative Policy Optimization (GRPO), which offers both high efficiency and effectiveness.

\subsubsection{From PPO to GRPO} Proximal Policy Optimization (PPO) \citep{schulman2017proximal} is an actor-critic RL algorithm that is widely used in the RL fine-tuning stage of LLMs \citep{ouyang2022training}. In particular, it optimizes LLMs by maximizing the following surrogate objective:

\begin{equation}
\footnotesize
    \mathcal{J}_{PPO}(\theta) = \mathbb{E}_{q \sim P(Q), o \sim \pi_{\theta_{old}}(O|q)} \left[ \frac{1}{|o|} \sum_{t=1}^{|o|} \min \left( \frac{\pi_\theta(o_{t} | q, o_{<t})}{\pi_{\theta_{old}}(o_{t} | q, o_{<t})} A_{t}, \ \text{clip} \left( \frac{\pi_\theta(o_{t} | q, o_{<t})}{\pi_{\theta_{old}}(o_{t} | q, o_{<t})},\ 1 - \varepsilon,\ 1 + \varepsilon \right)  A_{t} \right) \right],
\end{equation}

Here, $\pi_{\theta}$ denotes the updated policy model while $\pi_{\theta_{old}}$ represents the prior policy version; $q$ and $o$ correspond to question-response pairs, with questions drawn from a designated question set and outputs sampled using $\pi_{\theta_{old}}$. The hyper-parameter $\varepsilon$ is associated with clipping in PPO and is introduced to enhance training stability. $A_t$ refers to the computed advantage, which leverages the Generalized Advantage Estimation (GAE) framework \citep{gae}, utilizing subsequent rewards $\{r_{\ge t}\}$ and a learned value function $V_{\psi}$. Therefore, training the value function is an essential part of PPO alongside optimizing the policy. Additionally, to prevent excessive reward exploitation, it is common practice to include a token-wise KL divergence penalty with respect to a reference model within the reward computation at each generation step, as proposed by \citep{ouyang2022training}:

\begin{equation}
    r_{t} = r_\varphi(q, o_{\le t}) - \beta \log\frac{\pi_{\theta}(o_{t}|q, o_{<t})}{\pi_{ref}(o_{t}|q, o_{<t})},
\label{eq:PPO-reward}
\end{equation}

In this context, $r_\varphi$ denotes the reward model, $\pi_{ref}$ represents the reference model—typically corresponding to the initial SFT model—and $\beta$ serves as the scaling factor for the KL divergence penalty.

As the value function employed in PPO is typically another model of comparable size as the policy model, it brings a substantial memory and computational burden. Additionally, during RL training, the value function is treated as a baseline in the calculation of the advantage for variance reduction. While in the LLM context, usually only the last token is assigned a reward score by the reward model, which may complicate the training of a value function that is accurate at each token. To address this, as shown in Figure \ref{fig:grpo}, we propose Group Relative Policy Optimization (GRPO), which obviates the need for additional value function approximation as in PPO, and instead uses the average reward of multiple sampled outputs, produced in response to the same question, as the baseline. More specifically, for each question $q$, GRPO samples a group of outputs $\{o_1, o_2, \cdots, o_G\}$  from the old policy  $\pi_{\theta_{old}}$  and then optimizes the policy model by maximizing the following objective:

\begin{equation}
\footnotesize
\begin{split}
    \mathcal{J}_{GRPO}(\theta) &= \mathbb{E}{[q \sim P(Q), \{o_i\}_{i=1}^G \sim \pi_{\theta_{old}}(O|q)]}  \\
    & \frac{1}{G}\sum_{i=1}^G\frac{1}{|o_i|} \sum_{t=1}^{|o_i|} \left\{ \min \left[ \frac{\pi_\theta(o_{i,t} | q, o_{i,<t})}{\pi_{\theta_{old}}(o_{i,t} | q, o_{i,<t})} \hat{A}_{i,t}, \text{clip} \left( \frac{\pi_\theta(o_{i,t} | q, o_{i,<t})}{\pi_{\theta_{old}}(o_{i,t} | q, o_{i,<t})}, 1 - \varepsilon, 1 + \varepsilon \right)  \hat{A}_{i,t} \right] - \beta \mathbb{D}_{KL}\left[\pi_{\theta} || \pi_{ref}\right]\right\} ,
\end{split}
\label{eq:GRPO-obj}
\end{equation}

where $\varepsilon$ and $\beta$ are hyper-parameters, and $\hat{A}_{i,t}$  is the advantage calculated based on relative rewards of the outputs inside each group only, which will be detailed in the following subsections. The group relative way that GRPO leverages to calculate the advantages, aligns well with the comparative nature of rewards models,  as reward models are typically trained on datasets of comparisons between outputs on the same question. Also note that, instead of adding KL penalty in the reward, GRPO regularizes by directly adding the KL divergence between the trained policy and the reference policy to the loss, avoiding complicating the calculation of  $\hat{A}_{i,t}$. And different from the KL penalty term used in (\ref{eq:PPO-reward}), we estimate the KL divergence with the following unbiased estimator \citep{kl_approx}:

\begin{equation}
\small
    \mathbb{D}_{KL}\left[\pi_{\theta} || \pi_{ref}\right] = \frac{\pi_{ref}(o_{i,t}|q,o_{i,<t})}{\pi_{\theta}(o_{i,t}|q,o_{i,<t})}- \log\frac{\pi_{ref}(o_{i,t}|q,o_{i,<t})}{\pi_{\theta}(o_{i,t}|q,o_{i,<t})} - 1,
\end{equation}

which is guaranteed to be positive.

\begin{algorithm}[t]
  \small
  \caption{Iterative Group Relative Policy Optimization}
  \textbf{Input} initial policy model $\pi_{\theta_{\text{init}}}$; reward models $r_\varphi$; task prompts $\mathcal{D}$; 
  hyperparameters $\varepsilon$, $\beta$, $\mu$
  \begin{algorithmic}[1]
    \State Initialize policy model: $\pi_\theta \leftarrow \pi_{\theta_{\text{init}}}$
    \For{iteration = 1, \dots, I}
       \State Set reference model: $\pi_{ref} \leftarrow \pi_{\theta}$
      \For{step = 1, \dots, M}
      \State Draw a mini-batch $\mathcal{D}_b$ from $\mathcal{D}$
      \State Assign $\pi_{\theta_{old}} \leftarrow \pi_{\theta}$ 
      \State For each prompt $q \in \mathcal{D}_b$, generate $G$ candidate responses $\{o_i\}_{i=1}^G \sim \pi_{\theta_{old}} (\cdot \mid q) $
      \State For every candidate $o_i$, evaluate rewards $\{r_i\}_{i=1}^{G}$ using $r_{\varphi}$ 
      \State Calculate $\hat{A}_{i,t}$ for each token $t$ in $o_i$ with group-relative advantage computation.
      \For{each GRPO update = 1, \dots, $\mu$}
        \State Optimize $\pi_{\theta}$ by maximizing the GRPO objective function (Equation \ref{eq:GC-GRPO})
      \EndFor
    \EndFor \State Continue training $r_\varphi$ by leveraging experience replay.
    \EndFor 
  \end{algorithmic}
  \textbf{Output} $\pi_\theta$
  \label{alg:iter-grpo}
\end{algorithm}

\subsubsection{Outcome Supervision RL with GRPO} More specifically, given a question $q$, a set of responses $\{o_1, o_2, \cdots, o_G\}$ is generated by sampling from the pre-existing policy $\pi_{\theta_{old}}$. Each output is then evaluated by a reward model, producing a reward vector $\mathbf{r} = \{r_1, r_2, \cdots, r_G\}$ for the group. These reward values are standardized by removing the mean of the group and dividing by their standard deviation. In outcome supervision, the normalized reward assigned to an entire output $o_i$ is propagated uniformly to all tokens in that output, effectively using the normalized reward as the advantage estimate: $\hat{A}_{i, t} = \widetilde{r}_i = \frac{r_i- {\rm mean}(\mathbf{r})}{{\rm std}(\mathbf{r})}$. The policy parameters are then updated to maximize the objective as expressed in equation (\ref{eq:GRPO-obj}).

\subsubsection{Process Supervision RL with GRPO} Relying solely on outcome supervision, which issues rewards exclusively after the final output is generated, may prove inadequate and inefficient for guiding the policy in intricate mathematical problems. To address this, in line with the methodology in \cite{wang2023math}, we examine process supervision. This approach delivers feedback after every reasoning step, rather than just at the conclusion. Specifically, for a given question $q$ and a set of $G$ sampled completions $\{o_1, o_2, \cdots, o_G\}$, a process reward model is employed to evaluate each intermediate step, generating a corresponding reward structure: $\mathbf{R} = \{ \{r_1^{index(1)},\cdots,r_1^{index(K_1)}\}, \cdots,  \{r_G^{index(1)},\cdots,r_G^{index(K_G)}\} \}$, where $index(j)$ designates the position of the $j$-th step’s terminal token, and $K_i$ denotes the total steps in output $i$. These step-wise rewards are subsequently standardized using the mean and standard deviation, so that $\widetilde{r}_i^{index(j)} = \frac{r_i^{index(j)} - {\rm mean(\mathbf{R})}}{{\rm std(\mathbf{R})}}$.
Afterwards, the advantage estimate for each token is calculated by accumulating the normalized rewards assigned to all succeeding reasoning steps: $\hat{A}_{i, t} = \sum_{index(j) \ge t} \widetilde{r}_i^{index(j)}$.
This process supervision then updates the policy by maximizing the objective outlined in equation (\ref{eq:GRPO-obj}).

\subsubsection{Iterative RL with GRPO} During the advancement of reinforcement learning training, the previous reward model may become inadequate for effectively guiding the updated policy model. To address this, we investigate iterative RL using GRPO. As detailed in Algorithm \ref{alg:iter-grpo}, our iterative GRPO framework involves generating fresh datasets for reward model training by sampling outputs from the current policy model. We apply a replay mechanism, which retains 10\% of previous data, to repeatedly update the reward model. Subsequently, the reference model is synchronized with the latest policy model, and we proceed to further refine the policy model employing the newly updated reward model.

\subsection{Training and Evaluating DeepSeekMath-RL}

We implement reinforcement learning (RL) utilizing DeepSeekMath-Instruct 7B. The RL training corpus comprises chain-of-thought-format items associated with GSM8K and MATH, extracted from the SFT data and totaling approximately 144K questions. To specifically assess the effect of RL on benchmarks that lack related examples during RL training, we omit all other SFT questions. Our approach to constructing the reward model's training dataset mirrors that in \citep{wang2023math}. The initial reward model is trained using DeepSeekMath-Base 7B and adopts a learning rate of 2e-5. For GRPO, the policy model employs a learning rate of 1e-6, with the KL coefficient fixed at 0.04. For each question, $64$ completions are generated. We use a maximum sequence length of 1024 tokens, and the batch size is set to 1024. The policy model undergoes a single update after each exploration cycle. DeepSeekMath-RL 7B is assessed on the same benchmarks as DeepSeekMath-Instruct 7B. Here, GSM8K and MATH—with chain-of-thought reasoning—are categorized as in-domain tasks for DeepSeekMath-RL 7B, while all remaining benchmarks are treated as out-of-domain tasks.

Table \ref{tab:sft_rl_math} presents a comparative analysis of both open-source and closed-source models employing chain-of-thought and tool-integrated reasoning strategies across English and Chinese benchmarks. The results reveal the following insights: 1) Leveraging chain-of-thought reasoning, DeepSeekMath-RL 7B achieves accuracy rates of 88.2\% on GSM8K and 51.7\% on MATH. These outcomes exceed those achieved by every open-source model within the 7B to 70B range, in addition to most closed-source alternatives. 2) Notably, DeepSeekMath-RL 7B’s training exclusively involved chain-of-thought formatted instruction tuning data from GSM8K and MATH, with DeepSeekMath-Instruct 7B serving as the initialization point. Despite this narrow focus in its training set, DeepSeekMath-RL 7B demonstrates improved performance over DeepSeekMath-Instruct 7B on every metric evaluated, highlighting the substantial benefits provided by reinforcement learning.

\section{Discussion}
Here, we present the outcomes of our investigations in both pre-training and reinforcement learning experiments.

\subsection{Lessons Learnt in Pre-Training}

We begin by presenting insights from our pre-training process. Unless explicitly stated otherwise, the configurations described in Section \ref{sec:quality-policy} will be maintained. Additionally, it should be clarified that throughout this section, references to the DeepSeekMath Corpus pertain to the 89B-token version obtained from the second cycle of data collection.

\subsubsection{Code Training Benefits Mathematical Reasoning}

An often-cited but not conclusively established hypothesis proposes that learning to code enhances reasoning capabilities. In this work, we aim to partially address this claim, focusing specifically on mathematics: code training enhances models' proficiency in mathematical reasoning, regardless of whether tools are used.

To study how code training affects mathematical reasoning, we experimented with the following two-stage training and one-stage training settings:

\textbf{Two-Stage Training}

\begin{itemize}[topsep=0pt]
    \item \textbf{Code Training for 400B Tokens $\rightarrow$ Math Training for 150B Tokens}: We pretrain DeepSeek-LLM 1.3B using 400 billion tokens comprised of code, subsequently fine-tuning on 150 billion tokens containing mathematical content;
    \item \textbf{General Training for 400B Tokens $\rightarrow$ Math Training for 150B Tokens}: For comparison, we conduct a control setup where the initial pretraining involves 400 billion tokens randomly drawn from an extensive general-purpose dataset compiled by DeepSeek-AI, instead of code tokens, to examine the role of code-based pretraining versus general language pretraining on downstream mathematical reasoning tasks.
\end{itemize}

\textbf{One-Stage Training}

\begin{itemize}[topsep=0pt]
    \item \textbf{Math Training on 150B Tokens}: DeepSeek-LLM 1.3B is subjected to training solely on 150B tokens derived from mathematics data.
    \item \textbf{Joint Training with 400B Code Tokens and 150B Math Tokens}: Sequential math training after code pretraining leads to diminished coding abilities. We therefore explore if integrating code and math tokens into a unified training phase can concurrently foster mathematical reasoning skills and reduce the risk of catastrophic forgetting.
\end{itemize}

\paragraph{Results}
As shown in Table \ref{tab:code-math} and Table \ref{tab:code-math-general-eval}, the downstream outcomes for various training configurations are presented.

Code training enhances program-aided mathematical reasoning, regardless of whether it is implemented in a two-stage or one-stage training framework. As illustrated in Table \ref{tab:code-math}, when utilizing a two-stage approach, code-focused training on its own leads to substantial gains in solving both GSM8K and MATH tasks via Python. Introducing math-specific training in the second stage leads to additional progress. Notably, in a one-stage training paradigm, interleaving code and math tokens not only alleviates the problem of catastrophic forgetting often encountered with two-stage processes, but also improves coding skills (Table \ref{tab:code-math-general-eval}) and bolsters program-based mathematical reasoning (Table \ref{tab:code-math}).

Code training likewise enhances mathematical reasoning capabilities even in the absence of external tools. In the context of two-stage training, the first phase involving code yields measurable improvements. This phase further facilitates more effective math training in the next stage, culminating in optimal performance outcomes. Conversely, merging code and math tokens in a single-stage training paradigm appears to diminish mathematical reasoning abilities when tool support is not available. A possible explanation is that DeepSeek-LLM 1.3B, constrained by its relatively small size, is unable to adequately integrate both code and mathematical content when exposed to them together.

\subsubsection{ArXiv Papers Seem Ineffective in Improving Mathematical Reasoning}

ArXiv papers are commonly included as a component of math pre-training data \citep{minerva,gpt-f,llemma,mathpile}. However, detailed analysis regarding their impact on mathematical reasoning has not been extensively conducted. Perhaps counter-intuitively, according to our experiments, arXiv papers seem ineffective in improving mathematical reasoning. We experiment with models of different sizes, including DeepSeek-LLM 1.3B and DeepSeek-Coder-Base-v1.5 7B \citep{deepseek-coder}, using arXiv corpora that underwent varied processing pipelines:

\begin{itemize}[topsep=0pt]
    \item \textbf{MathPile} \citep{mathpile}: a collection containing 8.9 billion tokens, curated using a set of cleaning and filtering strategies, with scientific arXiv articles comprising more than 85\% of the dataset;
    \item \textbf{ArXiv-RedPajama} \citep{redpajama}: includes all available arXiv LaTeX source files, systematically stripped of preambles, annotations, macros, and reference lists, resulting in a dataset of 28.0 billion tokens.
\end{itemize}

In our study, we individually trained DeepSeek-LLM 1.3B on 150B tokens and DeepSeek-Coder-Base-v1.5 7B on 40B tokens using each arXiv dataset. The results indicate that utilizing arXiv papers alone does not enhance mathematical reasoning abilities. When the models are trained exclusively on arXiv data, neither model demonstrates significant gains; in fact, their performance either stagnates or declines across a range of mathematical evaluation sets with varying difficulty levels. These evaluation sets comprise quantitative reasoning benchmarks such as GSM8K and MATH (see Table \ref{tab:arxiv-cot}), multiple-choice tests including MMLU-STEM (see Table \ref{tab:arxiv-cot}), as well as formal mathematics tasks such as miniF2F (refer to Table \ref{tab:arxiv_atp}).

However, this conclusion has its limitations and should be taken with a grain of salt. We have not yet studied:

\begin{itemize}[topsep=0pt]
    \item How arXiv tokens influence mathematical tasks beyond the current scope of this study, for instance, the process of informalizing theorems, which entails translating formal statements or proofs into more informal language;
    \item The outcomes of integrating arXiv tokens with alternative forms of data;
    \item If scaling up to larger models would reveal additional advantages associated with arXiv papers.
\end{itemize}

Thus, further exploration is required, which we leave for future studies.

\subsection{Insights of Reinforcement Learning}
\subsubsection{Toward a Unified Framework} This section introduces a comprehensive framework that allows for the examination of a variety of training approaches, including SFT, RFT, DPO, PPO, and GRPO. We also perform experimental analyses to investigate the contributing elements of this unified framework. Broadly, the gradient concerning the parameter $\theta$ for any given training method can be expressed as:

\begin{equation}
    \nabla_{\theta}\mathcal{J}_{\textcolor{red}{\mathcal{A}}}(\theta) = \mathbb{E}[\underbrace{(q,o) \sim \textcolor{red}{\mathcal{D}}}_{Data \ Source}]\left( \frac{1}{|o|} \sum_{t=1}^{|o|}  \underbrace{GC_{{\mathcal{A}}}(q, o, t, \textcolor{red}{\pi_{{rf}}})}_{Gradient \ Coefficient}  \nabla_{\theta}\log \pi_{\theta}(o_t | q, o_{<t})\right).
\label{eq:objective}
\end{equation}

There exist three key components: 1) \textit{Data Source $\mathcal{D}$}, which determines the training data; 2) \textit{Reward Function $\pi_{{rf}}$}, which is the source of the training reward signal; 3) \textit{Algorithm $\mathcal{A}$}: which processes the training data and the reward signal to the gradient coefficient $GC$ that determines the magnitude of the penalty or reinforcement for the data. We analyze several representative methods based on such a unified paradigm:

\begin{itemize}
    \item  \textbf{Supervised Fine-tuning (SFT)}: SFT involves further training a pretrained model on a curated dataset selected by humans for fine-tuning.
    \item \textbf{Rejection Sampling Fine-tuning (RFT)}: RFT extends the SFT model's capabilities by conducting additional fine-tuning on outputs generated by the SFT model itself in response to SFT questions, retaining only those outputs deemed correct.
    \item \textbf{Direct Preference Optimization (DPO)}: DPO enhances the SFT model by additional fine-tuning on a set of expanded outputs derived from the SFT model, employing a pair-wise DPO loss for optimization.
    \item \textbf{Online Rejection Sampling Fine-tuning (Online RFT)}: Unlike the standard RFT process, Online RFT begins with a policy model initialized from the SFT model and continuously improves it via fine-tuning on outputs that are dynamically sampled from the evolving policy model.
    \item \textbf{PPO/GRPO}: PPO/GRPO leverages the SFT model for policy initialization and subsequently applies reinforcement learning, updating the model with outputs generated from the current policy in an online manner.
\end{itemize}

Table \ref{tab:data_policy} provides an overview of the elements comprising these approaches. For a comprehensive explanation of the derivation, see Appendix \ref{app:analysis-rl}.

\paragraph{Observation about Data Source} We categorize data sources into two main types: online sampling and offline sampling. Online sampling refers to collecting training data from the outputs generated by the current policy model during active exploration. In contrast, offline sampling involves utilizing data sampled from the preliminary SFT model's results. Methods such as RFT and DPO employ the offline approach, whereas Online RFT and GRPO utilize the online sampling methodology.

As illustrated in Figure \ref{fig:rl-analysis}, our results indicate that Online RFT consistently surpasses RFT on both evaluated benchmarks. Notably, performance between Online RFT and RFT is nearly indistinguishable during the initial training phase; however, as training progresses, Online RFT demonstrates a marked performance lead, underscoring the benefits of the online learning paradigm. This trend aligns with expectations: early on, the actor model and the SFT model are closely aligned, resulting in sampled data that reflects only minimal variations. Over time, as the actor diverges, the newly generated samples display more pronounced differences, making the advantages of continuous online data collection increasingly significant.

\paragraph{Observation about Gradient Coefficient} In the algorithm, the gradient coefficient, which influences model parameter updates, is determined by how the reward signal is processed. Within our experiments, the reward function is split into two types: `Rule' and `Model'. The `Rule' approach assesses response quality strictly by answer correctness, whereas the `Model' approach employs a trained reward model to assign a score to each response; notably, the training data for this model is itself derived from rule-based evaluations. As presented in Equations \ref{eq:GC-RFT} and \ref{eq:GC-GRPO}, a central distinction emerges between GRPO and Online RFT: only GRPO adapts its gradient coefficient in accordance with the actual reward output from the reward model. This capability enables GRPO to apply reinforcement or penalties to responses to varying extents, based on their respective reward scores. In comparison, Online RFT does not exhibit this adaptive behavior—it does not penalize incorrect answers at all, and it applies uniform positive reinforcement to all correct responses regardless of their individual quality.

As illustrated in Figure \ref{fig:rl-analysis}, GRPO outperforms online RFT, which underscores the effectiveness of modifying positive and negative gradient coefficients. Moreover, GRPO+PS achieves better results than GRPO+OS, demonstrating the advantage of employing step-aware, fine-grained gradient coefficients. We also investigate the impact of iterative RL by performing two consecutive iterations in our experiments. According to Figure \ref{fig:iter-rl}, iterative RL markedly enhances performance, with the most substantial gains observed during the initial iteration.

\subsubsection{Why RL Works?} In this study, we apply reinforcement learning to a selected subset of instruction tuning data, resulting in substantial performance improvements over the baseline instruction tuning model. To better understand the mechanism behind reinforcement learning's effectiveness, we analyze the Pass@K and Maj@K accuracies for both Instruct and RL models across two separate benchmarks. According to the results in Figure \ref{fig:maj-pass}, reinforcement learning brings a noticeable gain in Maj@K but leaves Pass@K relatively unchanged. This pattern suggests that reinforcement learning primarily strengthens the robustness of the model’s output distribution, that is, \textbf{the observed gains are more likely due to an increased probability of selecting a correct response among the TopK predictions, rather than an inherent advancement in the model’s core abilities.} Likewise, \citep{wang2023making} have observed a \textbf{misalignment problem} in reasoning tasks involving SFT models, and have demonstrated that leveraging various preference alignment methods \citep{yuan2023rrhf,song2023preference,wang2023making} can lead to enhanced reasoning performance in these models.

\subsubsection{How to Achieve More Effective RL?}
\label{sec:effec_rl}
Our findings indicate that RL is highly effective for tasks involving mathematical reasoning. Furthermore, we offer a comprehensive framework that clarifies the relationship among various prominent training strategies, situating each as a form of either direct or streamlined RL approaches. According to Equation \ref{eq:objective}, three fundamental elements can be identified: Data Source, Algorithm, and Reward Function. We outline several possible avenues for further exploration regarding these components.

\paragraph{Data Source} The data source serves as the foundational input for all training approaches. Within the framework of RL, we define the data source as comprising unlabeled questions alongside their associated outputs, which are generated by sampling from the policy model. For this study, our data consists exclusively of questions collected during the instruction tuning phase, and the output responses are produced using straightforward nucleus sampling. We suspect that limiting our data collection in this way may contribute to the observed improvement only in Maj@K performance within our RL framework. Looking ahead, we intend to evaluate the RL process using out-of-distribution question prompts and integrate \textbf{advanced sampling (decoding) strategies} informed by tree-search algorithms \citep{yao2023tree}. Furthermore, advancements in \textbf{efficient inference techniques} \citep{xia-etal-2023-speculative,leviathan2023fast,kwon2023efficient,xia2024unlocking}, which govern how effectively the policy model can explore, are anticipated to play a critical role in future enhancements.

\paragraph{Algorithms} The algorithms utilize both the dataset and the reward signal to compute the gradient coefficient used for model parameter updates. In accordance with Equation \ref{eq:objective}, contemporary approaches typically place complete \textbf{TRUST} in the reward function, adjusting the conditional probability of a token upward or downward depending on the signal received. Nonetheless, the reliability of the reward signal cannot always be guaranteed, particularly in highly complex tasks. As an illustration, even datasets like PRM800K \citep{lightman2023let}, despite having undergone rigorous annotation by proficient annotators, still exhibit an estimated 20\% rate of erroneous annotations\footnote{\url{https://github.com/openai/prm800k/issues/12\#issuecomment-1728491852}}. Given this limitation, our focus will be on examining reinforcement learning algorithms that maintain robustness in the presence of noisy reward signals. We anticipate that \textbf{WEAK-TO-STRONG} \citep{burns2023weak} alignment techniques have the potential to fundamentally reshape current learning algorithm paradigms.

\paragraph{Reward Function} The reward function serves as the primary source of the learning signal during training. In reinforcement learning (RL), this is most often implemented as a neural-based reward model. There are three key aspects regarding reward models that warrant significant attention: 1) \textbf{Strategies to improve reward model generalization.} It is crucial for the reward model to generalize robustly, particularly when faced with out-of-distribution prompts or sophisticated decoding scenarios; without sufficient generalization, RL risks only consolidating the LLM’s distribution rather than elevating its intrinsic performance; 2) \textbf{Capturing and utilizing reward model uncertainty.} Accounting for uncertainty in the reward model could provide a means of connecting underpowered reward models with more effective weak-to-strong learning frameworks; 3) \textbf{Scalable construction of precise process reward models} that yield granular training signals supporting the model’s reasoning procedures \citep{lightman2023let,wang2023math}.

\section{Conclusion, Limitation, and Future Work} In this work, we introduce DeepSeekMath, a model that surpasses all available open-source models on the MATH benchmark at the competition level, and achieves performance comparable to several closed-source systems. DeepSeekMath~is built upon DeepSeek-Coder-v1.5 7B, and is continually trained over 500B tokens, incorporating a substantial portion—120B tokens—of mathematical data extracted from Common Crawl. Our comprehensive ablation analysis reveals that web-sourced pages are a valuable resource for collecting high-quality mathematical content, while data from arXiv does not appear as advantageous as anticipated. Additionally, we propose Group Relative Policy Optimization (GRPO), an adaptation of Proximal Policy Optimization (PPO), which significantly boosts mathematical reasoning performance while reducing memory requirements. Empirical results demonstrate that GRPO remains effective even after DeepSeekMath-Instruct 7B attains high benchmark scores. Furthermore, we establish a coherent framework to analyze several related approaches and outline multiple promising avenues for improving reinforcement learning in future research.

While DeepSeekMath~demonstrates strong performance on quantitative reasoning tasks, its abilities in geometry and theorem proving are comparatively inferior to those of proprietary models. For example, during our preliminary evaluation, the model failed to address problems involving triangles and ellipses, which may reflect a bias in the types of data used during pre-training and fine-tuning stages. Furthermore, due to limitations in model size, DeepSeekMath~falls short of GPT-4 regarding few-shot learning capabilities; GPT-4 benefits significantly from few-shot prompts, whereas DeepSeekMath~exhibits nearly identical results under zero-shot and few-shot conditions. Moving forward, we plan to refine our engineered data selection process to build a more robust and higher-quality pre-training corpus. Moreover, we intend to investigate avenues described in Section \ref{sec:effec_rl} to realize more effective reinforcement learning strategies for LLMs.

\end{document}